\documentclass[11pt]{article}

\pdfinfoomitdate=1
\pdftrailerid{}
\pdfsuppressptexinfo=15

\usepackage[a4paper,margin=1in]{geometry}
\usepackage{amsmath,amssymb,amsthm,mathtools}
\usepackage{booktabs}
\usepackage{microtype}
\usepackage{xurl}
\usepackage[hidelinks]{hyperref}
\hypersetup{pdfcreator={},pdfproducer={}}
\usepackage[nameinlink,capitalise]{cleveref}

\newtheorem{theorem}{Theorem}[section]
\newtheorem{lemma}[theorem]{Lemma}
\newtheorem{proposition}[theorem]{Proposition}
\newtheorem{corollary}[theorem]{Corollary}
\theoremstyle{definition}
\newtheorem{definition}[theorem]{Definition}
\newtheorem{remark}[theorem]{Remark}

\newcommand{\R}{\mathbb{R}}
\newcommand{\Q}{\mathbb{Q}}
\newcommand{\cS}{\mathcal{S}}
\newcommand{\cZ}{\mathcal{Z}}
\newcommand{\Vol}{\operatorname{Vol}}
\newcommand{\rank}{\operatorname{rank}}
\newcommand{\tr}{\operatorname{tr}}
\newcommand{\im}{\operatorname{im}}
\newcommand{\aff}{\operatorname{aff}}
\newcommand{\spanop}{\operatorname{span}}
\newcommand{\Sym}{\operatorname{Sym}}

\title{Symbolic Infinitesimal Flexes for a Volume-Rigidity Conjecture in All Dimensions \(d\geq 7\)}
\author{Prepared with the EULER system}
\date{}

\begin{document}
\maketitle

\begin{abstract}
Cruickshank, Jackson, and Tanigawa conjectured that if \(\cS\) is a
simplicial \(2\)-sphere on \(\lceil d/2\rceil+2\) vertices and \(Z\)
is a disjoint set of \(d-3\) vertices, then the hypergraph of
\((d-2)\)-faces of the simplicial cone \(\cS*Z\) is volume rigid in
\(\R^d\).  We give a symbolic counterexample for every \(d\geq 7\).
Put \(m=d-4\) and choose one distinguished cone vertex \(z_0\).  For a
generic placement, let \(D\) be the \(d\times m\) matrix of cone-vertex
differences \(p(z_i)-p(z_0)\).  A nonzero symmetric matrix
\(K\in\Sym_m\) with zero diagonal, \(1^{\mathsf T}K1=0\), and
\(\tr(KD^{\mathsf T}D)=0\) gives a symmetric affine motion
\(u(v)=DKD^{\mathsf T}(p(v)-p(z_0))\).  The first two equations kill
the facet contributions from the cone simplex, and the trace equation
kills the hyperedges containing all cone vertices.  Hence every
squared \((d-2)\)-volume derivative vanishes, while the motion is not
Euclidean-trivial.  Therefore \(H_{d-2}(\cS*Z)\) is not volume rigid in
\(\R^d\) for any \(d\geq 7\).  Finite modular computations for
\(4\leq d\leq 9\) are reported only as cross-checks; the note does not
claim an exact generic rank or exact generic corank formula.
\end{abstract}

\section{Introduction}

Cruickshank, Jackson, and Tanigawa introduced a framework for volume
rigidity of uniform hypergraphs and simplicial complexes
\cite{CruickshankJacksonTanigawa2026}.  Among their proposed base cases
for higher-dimensional results is the following statement
\cite[Conjecture~22]{CruickshankJacksonTanigawa2026}: for \(d\geq 4\),
if \(\cS\) is a simplicial \(2\)-sphere on
\(\lceil d/2\rceil+2\) vertices and \(Z\) is a disjoint set of
\(d-3\) vertices, then the hypergraph \(H_{d-2}(\cS*Z)\) is volume
rigid in \(\R^d\).

This note gives a uniform symbolic obstruction for all dimensions
\(d\geq 7\).  The construction depends only on the way the
\((d-2)\)-faces of \(\cS*Z\) meet the cone simplex on \(Z\).  It is
therefore independent of the combinatorial type of the simplicial
\(2\)-sphere \(\cS\).

The lower dimensions are not re-proved here.  The source paper proves
the corresponding cases \(d=4,5,6\), and the theorem below shows that
the proposed family fails in every remaining dimension.  The result is
an infinitesimal generic non-rigidity statement.  It does not address
global uniqueness, nongeneric placements, exact generic rank, exact
generic corank, or the broader manifold conjecture discussed in
\cite{CruickshankJacksonTanigawa2026}.

\section{Definitions}
\label{sec:definitions}

\subsection{Simplicial cones and their hypergraphs}

For a simplicial complex \(\cS\) and a finite set \(Z\) disjoint from
\(V(\cS)\), define
\[
  \cS*Z=\{X\cup Z': X\in\cS,\ Z'\subseteq Z\}.
\]
For a simplicial complex \(\mathcal K\), the hypergraph \(H_k(\mathcal
K)\) has vertex set \(V(\mathcal K)\) and one hyperedge for every
\(k\)-face of \(\mathcal K\); hence every hyperedge has \(k+1\)
vertices.

Fix \(d\geq 7\).  Let \(\cS\) be a simplicial \(2\)-sphere on
\[
  s=\left\lceil\frac d2\right\rceil+2
\]
vertices, and let
\[
  Z=\{z_0,z_1,\ldots,z_m\},
  \qquad m=d-4,
\]
be disjoint from \(V(\cS)\).  Thus \(|Z|=m+1=d-3\).  The
\((d-1)\)-element hyperedges of \(H_{d-2}(\cS*Z)\) are exactly
\begin{equation}
  e\cup Z \quad (e\in\cS^{(1)}),
  \qquad
  t\cup\bigl(Z\setminus\{z_i\}\bigr)
  \quad (t\in\cS^{(2)},\ 0\leq i\leq m).
  \label{eq:hyperedge-families}
\end{equation}
Indeed, such a face contains either two sphere vertices and all
\(m+1=d-3\) cone vertices, or three sphere vertices and \(m=d-4\) cone
vertices.

\subsection{Squared-volume Jacobian}

Let \(p:V(H_{d-2}(\cS*Z))\to\R^d\) be a placement.  It is
\emph{generic} if its scalar coordinates are algebraically independent
over \(\Q\).  For a hyperedge
\(\Delta=\{v_0,v_1,\ldots,v_{d-2}\}\), let \(M_\Delta(p)\) be the
\((d-2)\times d\) matrix whose \(i\)-th row is
\((p(v_i)-p(v_0))^{\mathsf T}\), and put
\[
  g_\Delta(p)=\det\bigl(M_\Delta(p)M_\Delta(p)^{\mathsf T}\bigr).
\]
Then
\[
  g_\Delta(p)=((d-2)!)^2\Vol_{d-2}(p(\Delta))^2.
\]
Let \(J_{\cS,d}(p)\) be the squared-volume Jacobian with one row
\(\nabla g_\Delta(p)\) for each hyperedge \(\Delta\).  At a generic
placement this Jacobian has the same rank and right kernel as the
volume-rigidity matrix of
\cite{CruickshankJacksonTanigawa2026}: the Gram determinant differs
from squared volume by a fixed nonzero factor, and the row normalizing
factors used in that paper are nonzero at a generic placement.

The Euclidean-trivial infinitesimal motions are
\[
  u(v)=A p(v)+t,
  \qquad A^{\mathsf T}=-A,\quad t\in\R^d.
\]
For a generic placement with at least \(d+1\) vertices, these form a
\[
  d+\binom d2=\binom{d+1}{2}
\]
dimensional subspace of \(\ker J_{\cS,d}(p)\).  The generic placement
here has
\[
  |V(\cS)|+|Z|
  =\left\lceil\frac d2\right\rceil+d-1
  \geq d+1
\]
vertices, so affine spanning is generic.

\section{Main result}

\begin{theorem}
\label{thm:main}
Let \(d\geq 7\), let \(\cS\) be any simplicial \(2\)-sphere on
\(\lceil d/2\rceil+2\) vertices, let \(Z\) be a disjoint
\((d-3)\)-element set, and let
\[
  p:V(H_{d-2}(\cS*Z))\to\R^d
\]
be generic.  Then \(J_{\cS,d}(p)\) has a non-Euclidean infinitesimal
motion in its kernel.  Consequently \(H_{d-2}(\cS*Z)\) is not volume
rigid in \(\R^d\), and Conjecture~22 of
\cite{CruickshankJacksonTanigawa2026} is false for every \(d\geq 7\).
\end{theorem}

The remainder of this section constructs the motion and checks all
hyperedge derivatives.  The equivalence between generic volume rigidity
and generic infinitesimal volume rigidity is Proposition~6 of
\cite{CruickshankJacksonTanigawa2026}.

\subsection{A trace formula for squared volume}

For a linear subspace \(U\subseteq\R^d\), write \(P_U\) for its
orthogonal projector.

\begin{lemma}
\label{lem:trace}
Let \(\Delta\) be a nondegenerate \((d-2)\)-simplex with direction
space \(U\), and let \(B\) be a symmetric \(d\times d\) matrix.  Under
the affine velocity field \(u(q)=B(q-q_0)\), where \(q_0\) is fixed,
\[
  \left.\frac{d}{dt}\right|_{t=0}
  g_\Delta\bigl(q_0+(I+tB)(p-q_0)\bigr)
  =2g_\Delta(p)\tr(P_U B).
  \label{eq:trace-derivative}
\]
The same formula holds with \(g_\Delta\) replaced by squared
\((d-2)\)-volume.
\end{lemma}

\begin{proof}
Translations do not affect edge differences.  If \(M=M_\Delta(p)\),
then the transformed difference matrix is
\(M_t=M(I+tB)^{\mathsf T}\).  Since \(B=B^{\mathsf T}\),
\[
  \left.\frac{d}{dt}\right|_{t=0}M_tM_t^{\mathsf T}
  =2MBM^{\mathsf T}.
\]
Jacobi's determinant identity gives
\begin{align*}
  \left.\frac{d}{dt}\right|_{t=0}\det(M_tM_t^{\mathsf T})
  &=2\det(MM^{\mathsf T})
    \tr\!\left((MM^{\mathsf T})^{-1}MBM^{\mathsf T}\right)\\
  &=2g_\Delta(p)
    \tr\!\left(M^{\mathsf T}(MM^{\mathsf T})^{-1}M B\right).
\end{align*}
The matrix \(M^{\mathsf T}(MM^{\mathsf T})^{-1}M\) is \(P_U\), which
proves \eqref{eq:trace-derivative}.  Division by \(((d-2)!)^2\) gives
the squared-volume version.
\end{proof}

\subsection{The auxiliary matrix}

Translate notation so that \(p(z_0)\) is the origin, and define
\[
  D=\begin{bmatrix}d_1&d_2&\cdots&d_m\end{bmatrix}
   =\begin{bmatrix}
      p(z_1)-p(z_0)&p(z_2)-p(z_0)&\cdots&p(z_m)-p(z_0)
    \end{bmatrix},
  \qquad G=D^{\mathsf T}D.
\]
For generic \(p\), the \(d\times m\) matrix \(D\) has full column rank.

Let \(\mathcal K_G\) be the vector space of symmetric \(m\times m\)
matrices satisfying
\begin{equation}
  K_{ii}=0\quad(1\leq i\leq m),
  \qquad
  1^{\mathsf T}K1=0,
  \qquad
  \tr(KG)=0.
  \label{eq:K-space}
\end{equation}
The first condition leaves \(\binom m2\) off-diagonal variables.  The
last two conditions are two homogeneous linear equations in those
variables.  For generic \(G\) they are independent: \(1^{\mathsf T}K1\)
has coefficient vector proportional to the all-ones vector on
off-diagonal entries, while \(\tr(KG)\) has coefficient vector
proportional to the off-diagonal entries of \(G\), and these entries
are not generically all equal.  Hence
\[
  \dim \mathcal K_G=\binom m2-2.
\]
Since \(m=d-4\geq 3\), this dimension is positive.  Choose any nonzero
\(K\in\mathcal K_G\), and put
\begin{equation}
  B=DKD^{\mathsf T},
  \qquad
  u(v)=B\bigl(p(v)-p(z_0)\bigr).
  \label{eq:flex}
\end{equation}
Then \(B\) is symmetric and nonzero.  Indeed, if \(DKD^{\mathsf T}=0\)
and \(D\) has full column rank, then multiplying by a left inverse of
\(D\) and its transpose gives \(K=0\).

\begin{remark}
The dimension \(\binom m2-2\) is used here only to guarantee that a
nonzero \(K\) exists.  This note does not claim that it equals the
generic corank contribution.
\end{remark}

\subsection{Facet identities}

Let \(L=\im D\).  Inside \(L\), the cone simplex on \(Z\) has affine
coordinates \(0,e_1,\ldots,e_m\).  Its \(m+1\) facet normal covectors
are
\[
  e_1,\ldots,e_m,\quad 1=(1,\ldots,1)^{\mathsf T}.
\]
The zero diagonal and \(1^{\mathsf T}K1=0\) therefore give
\begin{equation}
  h^{\mathsf T}Kh=0
  \quad\text{for }h\in\{e_1,\ldots,e_m,1\}.
  \label{eq:facet-identities}
\end{equation}
The trace equation gives
\begin{equation}
  \tr(B)=\tr(DKD^{\mathsf T})=\tr(KG)=0.
  \label{eq:trace-B}
\end{equation}
Also \(B\) is supported on \(L\): it kills \(L^\perp\) and maps
\(\R^d\) into \(L\).

\subsection{Stationarity of all hyperedge volumes}

First consider a hyperedge \(\Delta=e\cup Z\).  Its direction space
has the orthogonal decomposition
\[
  U_\Delta=L\oplus R,
\]
where \(R\subseteq L^\perp\) is the span of the projections of the two
sphere vertices in \(e\).  Therefore
\begin{equation}
  \tr(P_{U_\Delta}B)=\tr(P_LB)+\tr(P_RB)=\tr(B)=0.
  \label{eq:edge-trace}
\end{equation}

Now let
\[
  \Delta=t\cup(Z\setminus\{z_\ell\}).
\]
Denote by \(L_\ell\) the \((m-1)\)-dimensional direction space of the
corresponding facet of the cone simplex on \(Z\).  Let \(h_\ell\) be
its normal covector in the \(D\)-coordinates.  Thus
\[
  h_\ell\in\{e_1,\ldots,e_m,1\}.
\]
The vector
\[
  n_\ell=DG^{-1}h_\ell
\]
is normal to \(L_\ell\) inside \(L\), because
\(D^{\mathsf T}n_\ell=h_\ell\).  From
\eqref{eq:facet-identities},
\begin{equation}
  n_\ell^{\mathsf T}Bn_\ell
  =h_\ell^{\mathsf T}Kh_\ell=0.
  \label{eq:normal-quadratic}
\end{equation}
Since \(L=L_\ell\oplus\spanop(n_\ell)\) orthogonally,
\eqref{eq:trace-B} and \eqref{eq:normal-quadratic} imply
\[
  \tr(P_{L_\ell}B)=0.
\]

The direction space of \(\Delta\) decomposes as
\[
  U_\Delta=L_\ell\oplus R_\ell
\]
for a subspace \(R_\ell\subseteq L_\ell^\perp\).  Moreover,
\[
  L_\ell^\perp=\spanop(n_\ell)\oplus L^\perp.
\]
For \(w=\lambda n_\ell+r\), with \(r\in L^\perp\), the facts
\(Br=0\) and \eqref{eq:normal-quadratic} give
\(w^{\mathsf T}Bw=0\).  Since \(B\) is symmetric, polarization shows
that the compression of \(B\) to \(L_\ell^\perp\) is zero.  Hence
\(\tr(P_{R_\ell}B)=0\), and altogether
\begin{equation}
  \tr(P_{U_\Delta}B)=0.
  \label{eq:triangle-trace}
\end{equation}

Equations \eqref{eq:edge-trace} and \eqref{eq:triangle-trace} cover
the two hyperedge families in \eqref{eq:hyperedge-families}.  By
\cref{lem:trace}, the velocity \(u\) in \eqref{eq:flex} lies in the
kernel of the generic squared-volume Jacobian.

\subsection{Nontriviality}

Suppose, for contradiction, that \(u\) is Euclidean-trivial.  Then
there are a skew-symmetric matrix \(A\) and a vector \(t\) such that
\[
  B(p(v)-p(z_0))=Ap(v)+t
\]
for every vertex \(v\).  Substitution of \(v=z_0\) gives
\(t=-Ap(z_0)\), so
\[
  B(p(v)-p(z_0))=A(p(v)-p(z_0))
\]
for every vertex \(v\).  A generic placement affinely spans \(\R^d\);
its differences from \(p(z_0)\) therefore span \(\R^d\), forcing
\(B=A\).  This is impossible because \(B\) is nonzero and symmetric
while \(A\) is skew-symmetric.  Thus the constructed kernel vector is
not Euclidean-trivial, proving \cref{thm:main}.

\section{The dimension-seven specialization}
\label{sec:d7}

When \(d=7\), one has \(m=3\), \(|Z|=4\), and
\(|V(\cS)|=6\).  Write
\[
  G=D^{\mathsf T}D,\qquad
  a=G_{12},\quad b=G_{13},\quad c=G_{23}.
\]
Then \(\mathcal K_G\) is one-dimensional and may be generated by
\[
  \xi=c-b,\qquad \eta=a-c,\qquad \zeta=b-a,
  \qquad
  K=\begin{pmatrix}
      0&\xi&\eta\\
      \xi&0&\zeta\\
      \eta&\zeta&0
    \end{pmatrix}.
\]
Indeed,
\[
  \xi+\eta+\zeta=0,\qquad
  a\xi+b\eta+c\zeta=0,
\]
which are exactly the equations \(1^{\mathsf T}K1=0\) and
\(\tr(KG)=0\).

Every six-vertex simplicial \(2\)-sphere has twelve edges and eight
triangles.  Thus \(H_5(\cS*Z)\) has
\[
  12+4\cdot 8=44
\]
hyperedges.  A generic ten-point placement in \(\R^7\) has \(70\)
velocity coordinates and \(28\) Euclidean-trivial motions, so rigidity
would require rank \(42\).  The construction above gives at least one
additional nontrivial kernel vector, hence the generic rank is at most
\(41\).

\section{Finite cross-checks and verification boundary}
\label{sec:verification}

The proof of \cref{thm:main} is the symbolic argument in the preceding
sections.  The following finite computations are included only as
independent cross-checks of the hyperedge construction and derivative
matrices.  They do not certify exact generic rank or exact generic
corank.

A direct modular Jacobian checker reconstructed the relevant
simplicial \(2\)-sphere types from the graph atlas, formed all
hyperedges in \eqref{eq:hyperedge-families}, and differentiated the
Gram determinants over three large primes and three random seeds.  The
observed ranks were stable in the following tests.

\begin{center}
\small
\begin{tabular}{rrrrr}
\toprule
\(d\) & sphere types tested & rows & columns & observed rank\\
\midrule
4 & 1 & 10 & 20 & 10\\
5 & 1 & 21 & 35 & 20\\
6 & 1 & 27 & 48 & 27\\
7 & 2 & 44 & 70 & 40\\
8 & 2 & 52 & 88 & 47\\
9 & 5 & 75 & 117 & 63\\
\bottomrule
\end{tabular}
\end{center}

A separate affine-flex-space checker solved the equations
\(\tr(P_U B)=0\) for symmetric affine matrices \(B\) at the same
specialized placements.  Its observed dimensions were \(0,0,0,2,5,9\)
for \(d=4,5,6,7,8,9\), respectively.  For \(d\geq 7\), this is
consistent with the symbolic construction and with additional affine
motions in those finite specializations.  The public theorem uses only
the existence of one such non-Euclidean motion for generic placements.

The \(d=7\) package also contains the older exact hyperedge tables for
the two six-vertex sphere types.  They are retained for reproducibility
and link stability.

\section{Relation to the source paper and scope}

The source paper proves the lower-dimensional range \(d=4,5,6\) for
the relevant base cases and states Conjecture~22 as a proposed family
for extending the method
\cite{CruickshankJacksonTanigawa2026}.  The theorem here disproves
that conjecture in every dimension \(d\geq 7\).  It does not settle the
source paper's broader Conjecture~2, since that extension involves
additional manifold-rigidity arguments beyond this base-case family.

This note makes no claim of public priority, of being the first
counterexample, or of establishing the current public status of the
all-dimensional conjecture.  Public novelty or priority is
\texttt{NOT\_ESTABLISHED}.

\appendix
\section{Dimension-seven hyperedge lists}
\label{app:hyperedges}

In the tables below, a string such as \texttt{012678} denotes the
six-element set \(\{0,1,2,6,7,8\}\).  Each table contains all \(44\)
hyperedges in the order used for the older exact \(d=7\) cross-check.

\subsection*{Stacked type}

\begin{center}
\small
\begin{tabular}{llll}
\toprule
\texttt{012678}&\texttt{012679}&\texttt{012689}&\texttt{012789}\\
\texttt{013678}&\texttt{013679}&\texttt{013689}&\texttt{013789}\\
\texttt{016789}&\texttt{023678}&\texttt{023679}&\texttt{023689}\\
\texttt{023789}&\texttt{026789}&\texttt{036789}&\texttt{124678}\\
\texttt{124679}&\texttt{124689}&\texttt{124789}&\texttt{126789}\\
\texttt{134678}&\texttt{134679}&\texttt{134689}&\texttt{134789}\\
\texttt{136789}&\texttt{146789}&\texttt{235678}&\texttt{235679}\\
\texttt{235689}&\texttt{235789}&\texttt{236789}&\texttt{245678}\\
\texttt{245679}&\texttt{245689}&\texttt{245789}&\texttt{246789}\\
\texttt{256789}&\texttt{345678}&\texttt{345679}&\texttt{345689}\\
\texttt{345789}&\texttt{346789}&\texttt{356789}&\texttt{456789}\\
\bottomrule
\end{tabular}
\end{center}

\subsection*{Octahedral type}

\begin{center}
\small
\begin{tabular}{llll}
\toprule
\texttt{012678}&\texttt{012679}&\texttt{012689}&\texttt{012789}\\
\texttt{013678}&\texttt{013679}&\texttt{013689}&\texttt{013789}\\
\texttt{016789}&\texttt{024678}&\texttt{024679}&\texttt{024689}\\
\texttt{024789}&\texttt{026789}&\texttt{034678}&\texttt{034679}\\
\texttt{034689}&\texttt{034789}&\texttt{036789}&\texttt{046789}\\
\texttt{125678}&\texttt{125679}&\texttt{125689}&\texttt{125789}\\
\texttt{126789}&\texttt{135678}&\texttt{135679}&\texttt{135689}\\
\texttt{135789}&\texttt{136789}&\texttt{156789}&\texttt{245678}\\
\texttt{245679}&\texttt{245689}&\texttt{245789}&\texttt{246789}\\
\texttt{256789}&\texttt{345678}&\texttt{345679}&\texttt{345689}\\
\texttt{345789}&\texttt{346789}&\texttt{356789}&\texttt{456789}\\
\bottomrule
\end{tabular}
\end{center}

\bibliographystyle{alpha}
\bibliography{references}

\end{document}
